%% LyX 2.4.4 created this file.  For more info, see https://www.lyx.org/.
%% Do not edit unless you really know what you are doing.
\documentclass[english]{article}
\usepackage[T1]{fontenc}
\usepackage[latin9]{inputenc}
\usepackage{amsmath}
\usepackage{amsthm}
\usepackage{amssymb}
\usepackage{graphicx}
\usepackage{geometry}
\geometry{verbose,tmargin=1in,bmargin=1in,lmargin=1in,rmargin=1in}

\makeatletter

%%%%%%%%%%%%%%%%%%%%%%%%%%%%%% LyX specific LaTeX commands.
%% A simple dot to overcome graphicx limitations
\newcommand{\lyxdot}{.}


%%%%%%%%%%%%%%%%%%%%%%%%%%%%%% User specified LaTeX commands.
\usepackage{fancyhdr}
\usepackage{lastpage}
\pagestyle{fancy}
\usepackage{enumitem}
\usepackage{tikz}
\usetikzlibrary{shapes,snakes}
\usepackage{amsmath,amssymb}
\usepackage{multicol}
% Added by lyx2lyx
\setlength{\parskip}{\smallskipamount}
\setlength{\parindent}{0pt}

\makeatother

\usepackage{babel}
\begin{document}
\lhead{MTH 331(A) Dr.\,James Ashe}
\chead{\textbf{Helpful Results}}
\cfoot{\thepage{}~of~\pageref{LastPage}}
%%
%%\setenumerate[0]{leftmargin=12pt,itemindent=0pt,labelwidth=0pt}
\renewcommand{\labelenumi}{\textbf{\arabic{enumi}.}}

\subsection*{Chapter 9}

\textbf{Squeeze Theorem for Sequences. }Let $\{a_{n}\}$, $\{b_{n}\}$,
and $\{c_{n}\}$ with $a_{n}\leq b_{n}\leq c_{n}$ for all integers
\emph{n }greater than some index $N$. If ${\displaystyle \lim_{n\rightarrow\infty}a_{n}}={\displaystyle \lim_{n\rightarrow\infty}b_{n}=L}$,
then $\lim_{n\rightarrow\infty}c_{n}=L$.\vfill

\textbf{Growth Rate of Sequences. }The following sequences are ordered
according to increasing growth rates as $n\rightarrow\infty$: $\{\ln n\}\ll\{n^{p}\}\ll\{n^{p}\ln n\}\ll\{n^{p+s}\}\ll\{b^{n}\}\ll\{n!\}\ll\{n^{n}\}$.
The ordering applies for positive real numbers $p,r$, and $s$ and
$b>1$.\vfill

\textbf{Geometric Series. }Let $r,a\in\mathbb{R}$. If $|r|<1$, then
${\displaystyle \sum_{k=0}^{\infty}}ar^{k}=\frac{a}{1-r}$. If $|r|\geq1$,
then the series diverges. \vfill

\textbf{Divergence Test. }If $\sum a_{k}$ converges, then ${\displaystyle \lim_{n\rightarrow\infty}a_{k}=0}.$
Equivalently, if ${\displaystyle \lim_{n\rightarrow\infty}a_{k}\neq0}$,
then the series diverges. \vfill

\textbf{Integral Test. }Suppose $f$ is a continuous, positive, decreasing
function for $x\geq1$ and let $a_{k}=f(k)$ for $k=1,2,3,\dots$.
Then 
\[
{\displaystyle \sum_{k=1}^{\infty}a_{k}}\mbox{ and \ensuremath{\int_{1}^{\infty}f(x)\,dx}}
\]
either both converge or both diverge. In the case of convergence,
the value of the integral is \emph{not}, in general, equal to the
value of the series. \vfill

\textbf{Ratio Test. }Let $\sum a_{k}$ be an infinite series with
positive terms and let $r={\displaystyle \lim_{n\rightarrow\infty}\frac{a_{k+1}}{a_{k}}}$.\vspace{-.25cm}
\begin{enumerate}
\item If $0\leq r<1$ then the series converges.
\item If $r>1$ (including $r=\infty$) then the series diverges.
\item If $r=1$ then the test is inconclusive.\vfill
\end{enumerate}
\textbf{Root Test. }Let $\sum a_{k}$ be an infinite series with positive
terms and let $\rho={\displaystyle \lim_{n\rightarrow\infty}\sqrt[k]{a_{k}}}$.\vspace{-.25cm}
\begin{enumerate}
\item If $0\leq\rho<1$ then the series converges.
\item If $\rho>1$ (including $r=\infty$) then the series diverges.
\item If $\rho=1$ then the test is inconclusive.\vfill
\end{enumerate}
\textbf{Comparison Test. }Let $\sum a_{k}$ and $\sum b_{k}$ be series
with positive terms.\vspace{-.25cm}
\begin{enumerate}
\item If $0<a_{k}\leq b_{k}$ and $\sum b_{k}$ converges, then $\sum a_{k}$
converges.
\item If $0<b_{k}\leq a_{k}$ and $\sum b_{k}$ diverges, then $\sum a_{k}$
diverges.\vfill
\end{enumerate}
\textbf{The Limit Comparison Test. }Let $\sum a_{k}$ and $\sum b_{k}$
be series with positive terms and let ${\displaystyle \lim_{k\rightarrow\infty}\frac{a_{k}}{b_{k}}=L}$.\vspace{-.25cm}
\begin{enumerate}
\item If $0<L<\infty$ then $\sum a_{k}$ and $\sum b_{k}$ either both
converge or diverge.
\item If $L=0$ and $\sum b_{k}$ converges, then $\sum a_{k}$ converges. 
\item If $L=\infty$ and $\sum b_{k}$diverges, then $\sum a_{k}$ diverges. 
\end{enumerate}

\subsection*{Chapter 10}

\textbf{Taylor's Polynomials. }Let $f$ be a function with $f',f'',\dots f^{(n)}$
defined at $a$. The \textbf{\emph{n}}\textbf{th-order Taylor Polynomial
}for $f$ with its \textbf{center }at $a$ denoted $p_{n}$, has the
property that it matches $f$ in value, slope, and all derivatives
up to the \emph{n}th derivative at $a$; that is,
\[
p_{n}(a)=f(a)\mbox{ and }p'_{n}(a)=f'(a),\dots,p_{n}^{(n)}(a)=f^{(n)}(a)\mbox{.}
\]

The \emph{n}th-order Taylor polynomial centered at $a$ is 
\[
p_{n}(x)=f(a)+f'(a)(x-a)+\frac{f''(a)}{2!}(x-a)^{2}+\dots+\frac{f^{(n)}(a)}{n!}(x-a)^{n}\mbox{.}
\]

More compactly, 
\[
p_{n}(x)={\displaystyle \sum_{k=0}^{n}\frac{f^{(k)}(a)}{k!}(x-a)^{k}\mbox{ for }k=0,1,2,\dots,n\mbox{.}}
\]
\vfill

\subsection*{Chapter 11}

\textbf{Slope of a Tangent Line. }Let $f$ be differentiable function
at $\theta_{0}$. The slope of the line tangent to the curve $r=f(\theta)$
at the point $\left(f(\theta),\theta_{0}\right)$ is 
\[
\frac{dy}{dx}=\frac{f'(\theta_{0})\sin\theta_{0}+f(\theta_{0})\cos\theta_{0}}{f'(\theta_{0})\cos\theta_{0}-f(\theta_{0})\sin\theta_{0}}\mbox{,}
\]
 provided the denominator is nonzero at the point. At angles $\theta_{0}$
for which $f(\theta_{0})=0$ and $f'(\theta_{0})\neq0$, the tangent
line is $\theta=\theta_{0}$ with slope $\tan\theta_{0}$.\vfill

\textbf{Area of Regions in Polar Coordinates. }Let \textbf{$R$ }be
the region bounded by the graph of $r=f(\theta)$ and the origin,
between $\theta=\alpha$ and $\theta=\beta$, where $f$ is continuous
and $f(\theta)\geq g(\theta)\geq0$ on $[\alpha,\beta]$. The area
of $R$ is 
\[
{\displaystyle \int_{\alpha}^{\beta}\frac{1}{2}f(\theta)^{2}d\theta}\mbox{.}
\]
\vfill

\subsection*{12.3 Dot Products}

\textbf{Dot Product. }Given two vectors $\mathbf{u}=\left\langle u_{1},u_{2},u_{3}\right\rangle $
and $\mathbf{v}=\left\langle v_{1},v_{2},v_{3}\right\rangle $, 
\begin{eqnarray*}
\mathbf{u\cdot v} & = & u_{1}v_{1}+u_{2}v_{2}+u_{3}v_{3}\\
 & = & |\mathbf{u}||\mathbf{v}|\cos\theta
\end{eqnarray*}
\vfill

\textbf{Orthogonal Projection of u onto v},\textbf{ 
\begin{eqnarray*}
\mbox{proj}_{\mathbf{v}}\mathbf{u} & = & \left|\mathbf{u}\right|\cos\theta\left(\frac{\mathbf{v}}{\left|\mathbf{v}\right|}\right)\\
 & = & \mbox{scal}_{\mathbf{v}}\mathbf{u}\left(\frac{\mathbf{v}}{\left|\mathbf{v}\right|}\right)\\
 & = & \left(\frac{\mathbf{u\cdot v}}{\mathbf{v\cdot v}}\right)\mathbf{v\mbox{.}}
\end{eqnarray*}
\vfill}

\textbf{Scalar component of u in the direction of v},
\[
\mbox{scal}_{\mathbf{v}}\mathbf{u}=\left|\mathbf{u}\right|\cos\theta=\frac{\mathbf{u\cdot v}}{\left|\mathbf{v}\right|}\mbox{.}
\]
\vfill

\textbf{Work. }Let \textbf{F }be a constant force applied to object,
producing a displacement \textbf{d}. If the angle between \textbf{F
}and \textbf{d} is $\theta$, then the \textbf{work }done by the force
is 
\[
W=\left|\mathbf{F}\right|\left|\mathbf{d}\right|\cos\theta=\mathbf{F\cdot d}.
\]


\subsection*{12.4 Cross Products}

\textbf{Cross product}, of two nonzero vectors \textbf{u} and \textbf{v
}in $\mathbb{R}^{3}$, 
\[
\left|\mathbf{u\times v}\right|=\left|\mathbf{u}\right|\left|\mathbf{v}\right|\sin\theta
\]
\[
\left|\tau\right|=\left|\mathbf{r}\right|\left|\mathbf{F}\right|\sin\theta
\]
\[
\mathbf{u\times v}=\left|\begin{array}{ccc}
\mathbf{i} & \mathbf{j} & \mathbf{k}\\
u_{1} & u_{2} & u_{3}\\
v_{1} & v_{2} & v_{3}
\end{array}\right|=\left|\begin{array}{cc}
u_{2} & u_{3}\\
v_{2} & v_{3}
\end{array}\right|\mathbf{i}-\left|\begin{array}{cc}
u_{1} & u_{3}\\
v_{1} & v_{3}
\end{array}\right|\mathbf{j}+\left|\begin{array}{cc}
u_{1} & u_{2}\\
v_{1} & v_{2}
\end{array}\right|\mathbf{k}\mbox{.}
\]

\vfill\newpage

\subsection*{13 Functions of Several Variables}

\textbf{Chain Rule (Two Independent Variables). }Let $z$ be a differentiable
function of $x$ and $y,$ where $x$ and $y$ are differentiable
functions of $s$ and $t$. Then 
\[
\frac{dz}{dt}=\frac{dz}{dx}\frac{dx}{dt}+\frac{dz}{dy}\frac{dy}{dt}\mbox{.}
\]
\vfill

\textbf{Gradient (Two Dimensions). }Let $f$ be differentiable at
the point $(x,y)$. The \textbf{gradient} of $f$ at $(x,y)$ is the
vector-valued function 
\[
\nabla f(x,y)=\left\langle f_{x}(x,y),f_{y}(x,y)\right\rangle =f_{x}(x,y)\mathbf{i}+f_{y}(x,y)\mathbf{j}\mbox{.}
\]
\vfill

\textbf{Directional Derivative. }Let $f$ be differentiable at $(a,b)$
and let $\mathbf{u}=\left\langle u_{1},u_{2}\right\rangle $ be a
unit vector in the \emph{xy-}plane, the \textbf{directional derivative
of $f$ at $\left(a,b\right)$ in the direction of u }is 
\[
D_{\mathbf{u}}f(a,b)=\left\langle f_{x}(a,b),f_{y}(a,b)\right\rangle \cdot\left\langle u_{1},u_{2}\right\rangle 
\]

\[
D_{\mathbf{u}}f(a,b)=\nabla f(a,b)\cdot\mathbf{u}\mbox{.}
\]
\vfill

\textbf{Second Derivative Test. }Suppose that the second partial derivatives
of $f$ are continuous throughout an open disk centered at the point
$(a,b)$, where $f_{x}(a,b)=f_{y}(a,b)=0$. Let $D(x,y)=f_{xx}f_{yy}-f_{xy}^{2}$.
\begin{enumerate}
\item If $D(a,b)>0$ and $f_{xx}(a,b)<0$, then $f$ has a local maximum
value at $(a,b)$.
\item If $D(a,b)>0$ and $f_{xx}(a,b)>0$, then $f$ has a local minimum
value at $(a,b)$.
\item If $D(a,b)<0$, then $f$ has a saddle point at $(a,b)$.
\item If $D(a,b)=0$, then the test is inconclusive.\vfill
\end{enumerate}
\begin{multicols}{2}

\subsection*{Helpful formulas}
\begin{enumerate}
\item $\frac{d}{dx}\left(x^{n}\right)=nx^{n-1}$
\item $\frac{d}{dx}\left[f\left(g(x)\right)\right]=f'\left(g(x)\right)\cdot g'(x)$
(Chain rule)
\item $\frac{d}{dx}\left(f(x)g(x)\right)=f'(x)g(x)+f(x)g'(x)$
\item $\frac{d}{dx}e^{x}=e^{x}$
\item $\frac{d}{dx}\sin ax=a\cos ax$
\item $\frac{d}{dx}\cos ax=-a\sin ax$
\item $\int x^{n}\,dx=\frac{1}{n+1}x^{n+1}+C$; $n\neq1$ 
\item $\int\cos ax\,dx=\frac{1}{a}\sin ax+C$
\item $\int\sin ax\,dx=-\frac{1}{a}\cos ax+C$
\item $\int\cos^{2}x\,dx=\frac{x}{2}+\frac{\sin2x}{4}+C$
\item $\int\sin^{2}x\,dx=\frac{x}{2}-\frac{\sin2x}{4}+C$%\vfill
\columnbreak
\end{enumerate}
\includegraphics[width=2.5in]{slides_Calc3/images/Unit_circle_angles_color.svg.png}

\end{multicols}
\end{document}
